\chapter{Overview}
\label{chap:overview}

% archon:covers MR4433080HitchinFibrationsAbelianSurfacesAndThePwConjecture.lean MR4433080HitchinFibrationsAbelianSurfacesAndThePwConjecture/Basic.lean

% Chapter-local macros (to be promoted to blueprint/src/macros/common.tex)
\providecommand{\MDol}{\mathcal{M}_{\mathrm{Dol}}}
\providecommand{\MB}{\mathcal{M}_B}
\providecommand{\QQ}{\mathbb{Q}}
\providecommand{\ZZ}{\mathbb{Z}}
\providecommand{\HodgeTate}[2]{{}^{#1}\!\mathrm{Hdg}^{#2}}
\providecommand{\ch}{\operatorname{ch}}
\providecommand{\supp}{\operatorname{supp}}
\providecommand{\rk}{\operatorname{rank}}

This chapter covers the P=W conjecture for Hitchin fibrations, as proved by
de~Cataldo, Maulik, and Shen for genus-\(2\) curves (arbitrary rank) and, in higher
genus, for the subalgebra generated by even tautological classes.  The main
theorems identify the perverse filtration on \(H^*(\MDol,\QQ)\) with the weight
filtration on \(H^*(\MB,\QQ)\) under the non-abelian Hodge diffeomorphism.  The
key technique is to study analogous filtrations on moduli spaces of sheaves on
abelian surfaces and transfer the results to the Hitchin setting via a
degeneration argument.  Mathematical infrastructure not yet in Mathlib (perverse
sheaves, mixed Hodge structures, non-abelian Hodge theory) is axiomatized in
Lean.

% The section/subsection structure below mirrors the reference paper
% (de~Cataldo--Maulik--Shen) in its own order: each blueprint item is grouped
% under the paper section/subsection from which it originates (per its
% % SOURCE: provenance comment).

\section{Introduction}

\subsection{The P=W conjecture}

% ============================================================
%  Definitions 3--7: The P=W setup
% ============================================================

% SOURCE: de~Cataldo--Maulik--Shen, Section~0.2
%   (read from references/MR4433080-hitchin-pw.tex)
% SOURCE QUOTE: The first moduli space $\CM_{\mathrm{Dol}}$ parametrizes stable
% Higgs bundles on $C$
% $(\CE, \theta), \quad \theta: \CE \to \CE \otimes \Omega_C$
% with $\mathrm{rank}(\CE)=r$ and $\chi(\CE) =\chi$.
\begin{definition}[Dolbeault moduli space]
  \label{def:HitchinModuliDol}
  \lean{HitchinPW.HitchinModuliDol}
  \uses{}
  \textit{Source: de~Cataldo--Maulik--Shen, Section~0.2.}
  Let \(C\) be a smooth projective curve of genus \(g \geq 2\), and fix
  integers \(r \geq 1\) and \(\chi\) with \(\gcd(r,\chi) = 1\).  The
  \emph{Dolbeault moduli space} \(\MDol\) parametrizes \(S\)-equivalence classes
  of stable Higgs bundles
  \[
    (\mathcal{E},\theta), \quad \theta \colon \mathcal{E} \to \mathcal{E} \otimes \Omega_C,
  \]
  with \(\rk(\mathcal{E}) = r\) and \(\chi(\mathcal{E}) = \chi\).
  It is a nonsingular quasi-projective variety.
\end{definition}

% SOURCE: de~Cataldo--Maulik--Shen, Section~0.2, equation~(0.2)
%   (read from references/MR4433080-hitchin-pw.tex)
% SOURCE QUOTE: The variety $\CM_{\mathrm{Dol}}$ admits a projective morphism
% with connected fibers
% $h: \CM_{\mathrm{Dol}} \rightarrow \Lambda = \bigoplus_{i=1}^n H^0(C, \Omega_C^{\otimes i})$
% sending $(\CE, \theta) \in \CM_{\mathrm{Dol}}$ to the characteristic polynomial
% $\mathrm{char}(\theta) \in \Lambda$.
\begin{definition}[Hitchin fibration]
  \label{def:HitchinFibration}
  \lean{HitchinPW.HitchinFibration}
  \uses{def:HitchinModuliDol}
  \textit{Source: de~Cataldo--Maulik--Shen, Section~0.2, equation~(0.2).}
  The \emph{Hitchin fibration} is the proper morphism with connected fibers
  \[
    h \colon \MDol \longrightarrow \Lambda
      = \bigoplus_{i=1}^r H^0(C, \Omega_C^{\otimes i}),
  \]
  sending \((\mathcal{E},\theta)\) to the characteristic polynomial
  \(\operatorname{char}(\theta) \in \Lambda\).  It is Lagrangian with respect
  to the canonical holomorphic symplectic form on \(\MDol\).
\end{definition}

% SOURCE: de~Cataldo--Maulik--Shen, Section~0.2
%   (read from references/MR4433080-hitchin-pw.tex)
% SOURCE QUOTE: The second moduli space is the (twisted) character variety
% $\CM_B$.  It can be described (see \cite{HT1}, or \cite{Shende} for a an
% alternative description) as parametrizing isomorphism classes of irreducible
% local systems
% $\rho: \pi_1(C \setminus \{p\}) \to \mathrm{GL}_r$
% where $\rho$ sends a loop around a chosen point $p$ to $\xi^\chi_r I_r \in
% \mathrm{GL}_r$, with $\xi_r$ a primitive root of unity.
\begin{definition}[Character variety]
  \label{def:CharacterVariety}
  \lean{HitchinPW.CharacterVariety}
  \uses{}
  \textit{Source: de~Cataldo--Maulik--Shen, Section~0.2.}
  The (twisted) \emph{character variety} \(\MB\) parametrizes isomorphism
  classes of irreducible representations
  \[
    \rho \colon \pi_1(C \setminus \{p\}) \to \mathrm{GL}_r
  \]
  sending a loop around the basepoint \(p \in C\) to \(\xi_r^\chi I_r\), where
  \(\xi_r\) is a fixed primitive \(r\)-th root of unity.  The variety \(\MB\)
  is affine.
\end{definition}

% SOURCE: de~Cataldo--Maulik--Shen, Section~0.2
%   (read from references/MR4433080-hitchin-pw.tex)
% SOURCE QUOTE: In \cite{Simp}, Simpson constructed a diffeomorphism between
% the (un-twisted) moduli spaces $\CM_{\mathrm{Dol}}$ and $\CM_B$, called the
% non-abelian Hodge theory; see \cite{HT1} for the twisted case.
% ... it predicts that the perverse filtration of $\CM_{\mathrm{Dol}}$ with
% respect to the Hitchin fibration (\ref{Hitchin_fib}) matches the weight
% filtration of the mixed Hodge structure on $\mathcal{M}_B$ under the
% identification $H^\ast(\CM_{\mathrm{Dol}}, \BQ) = H^\ast(\CM_B, \BQ)$
% induced by Simpson's (twisted) non-abelian Hodge theory.
\begin{definition}[Non-abelian Hodge diffeomorphism]
  \label{def:NonabelianHodgeDiffeo}
  \lean{HitchinPW.NonabelianHodgeDiffeo}
  \uses{def:HitchinModuliDol, def:CharacterVariety}
  \textit{Source: de~Cataldo--Maulik--Shen, Section~0.2.}
  Simpson's \emph{non-abelian Hodge diffeomorphism} is a diffeomorphism
  \(\MDol \xrightarrow{\sim} \MB\) that induces an isomorphism of graded
  \(\QQ\)-algebras
  \[
    H^*(\MDol, \QQ) \cong H^*(\MB, \QQ).
  \]
  This diffeomorphism is axiomatized (not proved) in the Lean formalization.
\end{definition}

% SOURCE: de~Cataldo--Maulik--Shen, Section~0.2
%   (read from references/MR4433080-hitchin-pw.tex)
% SOURCE QUOTE: it predicts that the perverse filtration of
% $\CM_{\mathrm{Dol}}$ with respect to the Hitchin fibration
% (\ref{Hitchin_fib}) matches the weight filtration of the mixed Hodge
% structure on $\mathcal{M}_B$ under the identification
% $H^\ast(\CM_{\mathrm{Dol}}, \BQ) = H^\ast(\CM_B, \BQ)$
% induced by Simpson's (twisted) non-abelian Hodge theory.
\begin{definition}[Weight filtration]
  \label{def:WeightFiltration}
  \lean{HitchinPW.WeightFiltration}
  \uses{def:CharacterVariety}
  \textit{Source: de~Cataldo--Maulik--Shen, Section~0.2.}
  The \emph{weight filtration} \(W_k H^d(\MB, \QQ)\) is the filtration on
  the rational cohomology of the character variety \(\MB\) arising from the
  mixed Hodge structure on \(\MB\).  It is axiomatized in the Lean
  formalization.
\end{definition}

% SOURCE: de~Cataldo--Maulik--Shen, Section~0.2, Theorem~0.1
%   (read from references/MR4433080-hitchin-pw.tex)
% SOURCE QUOTE: \begin{thm}\label{genus_2}
% The P=W
% Conjecture \ref{P=W_conj} holds when $C$ has genus $g=2$.
% \end{thm}
% [Conjecture P=W stated as:]
% $P_kH^\ast(\CM_{\mathrm{Dol}}, \BQ) = W_{2k}H^\ast(\CM_{B}, \BQ)=
%   W_{2k+1}H^\ast(\CM_{B}, \BQ), \quad k \geq 0.$
\begin{theorem}[P=W for genus~2]
  \label{thm:PW_genus2}
  \lean{HitchinPW.PW_genus2}
  \uses{def:PervFiltration, def:WeightFiltration, def:NonabelianHodgeDiffeo,
        def:HitchinFibration}
  \textit{Source: de~Cataldo--Maulik--Shen, Theorem~0.1.}
  Let \(C\) be a smooth projective curve of genus \(g = 2\), and let
  \(r \geq 1\), \(\gcd(r,\chi) = 1\).  Then
  \[
    P_k H^*(\MDol, \QQ) = W_{2k} H^*(\MB, \QQ) = W_{2k+1} H^*(\MB, \QQ)
    \quad \forall\, k \geq 0,
  \]
  under the identification \(H^*(\MDol,\QQ) = H^*(\MB,\QQ)\) given by the
  non-abelian Hodge diffeomorphism (see \cref{def:NonabelianHodgeDiffeo}).
\end{theorem}

% SOURCE QUOTE PROOF: [Section 4.6] Since the perverse filtration with
% respect to the Hitchin fibration is locally constant when we deform the
% curve, it suffices to prove the theorem for a curve embedded in an abelian
% surface $j_C: C \hookrightarrow A$.  We then have the specialization
% morphism $\mathrm{sp}^!: H^*(\CM_{\beta,A}, \BC) \to
% H^*(\CM_{\mathrm{Dol}}, \BC)$, and obtain: (i) $c(\gamma,k) \in
% \widetilde{G}_k H^*(\CM_{\mathrm{Dol}}, \BQ)$ for all $\gamma \in
% \mathrm{Im}(j_C^*)$; (ii) the restriction of the splitting to the
% subalgebra generated by these classes is multiplicative.  ``We first
% prove Theorem \ref{genus_2}. The Abel-Jacobi morphism embeds a genus 2
% curve $C$ into its Jacobian $j_C: C \hookrightarrow \mathrm{Jac}(C) = A$.
% Hence the restriction morphism $j_C^*$ is surjective, and Theorem
% \ref{genus_2} follows from (i) and (ii) above.''
\begin{proof}
  \uses{thm:taut_abelian_splitting, lem:taut_weight, thm:oddTautPerversity}
  The proof is the specialization argument of Section~4 of the paper.  Since
  the perverse filtration of the Hitchin fibration is locally constant under
  deformations of the curve, it suffices to treat a curve \(C\) that embeds
  into an abelian surface \(j_C \colon C \hookrightarrow A\).  One then has
  the specialization morphism
  \(\mathrm{sp}^! \colon H^*(\mathcal{M}_{\beta,A},\mathbb{C}) \to
  H^*(\MDol,\mathbb{C})\), which transports the abelian-surface splitting
  (\cref{thm:taut_abelian_splitting}) to a splitting
  \(\widetilde{G}_* H^*(\MDol,\QQ)\) satisfying:
  \begin{enumerate}
    \item[(i)] \(c(\gamma,k) \in \widetilde{G}_k H^*(\MDol,\QQ)\) for every
      \(\gamma \in \operatorname{Im}(j_C^*) \subset H^*(C,\QQ)\); and
    \item[(ii)] the restriction of the splitting to the subalgebra generated
      by the classes in (i) is multiplicative.
  \end{enumerate}
  For genus \(g = 2\), the Abel--Jacobi morphism embeds \(C\) into its
  Jacobian \(j_C \colon C \hookrightarrow \operatorname{Jac}(C) = A\), an
  abelian surface, so the restriction \(j_C^*\) is \emph{surjective} onto
  \(H^*(C,\QQ)\).  Hence (i) and (ii) hold for \emph{all} tautological
  generators \(c(\gamma,k)\) --- including the odd ones, whose perversity is
  the content of \cref{thm:oddTautPerversity} --- so the multiplicative
  splitting \(\widetilde{G}_*\) covers all of \(H^*(\MDol,\QQ)\).  Combined
  with the weight identity \cref{lem:taut_weight}
  (\(c(\gamma,k) \in {}^k\!\mathrm{Hdg}(\MB) \subset W_{2k}\)) this matches
  the perverse splitting with the weight decomposition.  The full P=W
  equality \(P_k = W_{2k} = W_{2k+1}\) then follows by the Curious Hard
  Lefschetz theorem of Mellit (Mellit's Theorem~1.5.3, applied as in
  \cref{thm:PW_iff_multiplicativity}).  Markman's monodromy operators,
  Mellit's theorem, and the specialization morphism are axiomatized in the
  Lean formalization; the full argument is in Section~4 of the paper.
\end{proof}

\subsection{Tautological classes}

% ============================================================
%  Definitions 8--12: Tautological classes
% ============================================================

% SOURCE: de~Cataldo--Maulik--Shen, Section~0.3
%   (read from references/MR4433080-hitchin-pw.tex)
% SOURCE QUOTE: We say that a triple
% $(\CU^\alpha, \theta) = (\CU, \theta, \alpha)$
% is a {\em twisted universal family} over $C \times \CM_{\mathrm{Dol}}$,
% if $(\CU, \theta)$ is a universal family and
% $\alpha = p_C^\ast \alpha_1 +p_\CM^\ast \alpha_2 \in H^2(C\times
%   \CM_{\mathrm{Dol}}, \BQ),$
% with $\alpha_1 \in H^2(C, \BQ)$ and $\alpha_2 \in H^2(\CM_{\mathrm{Dol}},
% \BQ)$.
\begin{definition}[Twisted universal family]
  \label{def:TwistedUniversalFamily}
  \lean{HitchinPW.TwistedUniversalFamily}
  \uses{def:HitchinModuliDol}
  \textit{Source: de~Cataldo--Maulik--Shen, Section~0.3.}
  A \emph{twisted universal family} over \(C \times \MDol\) is a triple
  \((\mathcal{U}^\alpha, \theta) = (\mathcal{U}, \theta, \alpha)\), where
  \((\mathcal{U},\theta)\) is a universal family of stable Higgs bundles and
  \[
    \alpha = p_C^* \alpha_1 + p_{\mathcal{M}}^* \alpha_2
      \in H^2(C \times \MDol, \QQ),
  \]
  with \(\alpha_1 \in H^2(C, \QQ)\) and
  \(\alpha_2 \in H^2(\MDol, \QQ)\).
\end{definition}

% SOURCE: de~Cataldo--Maulik--Shen, Section~0.3
%   (read from references/MR4433080-hitchin-pw.tex)
% SOURCE QUOTE: For a twisted universal family $(\CU^\alpha, \theta)$, we
% define the {\em twisted Chern character} $\mathrm{ch}^\alpha(\CU)$ as
% $\mathrm{ch}^\alpha(\CU) = \mathrm{ch}(\CU) \cup \mathrm{exp}(\alpha)
%   \in H^*(C\times \CM_{\mathrm{Dol}}, \BQ),$
% and we denote by $\mathrm{ch}_k^\alpha(\CU) \in H^{2k}(C\times
%   \CM_{\mathrm{Dol}},\BQ)$ its degree $k$ part.
\begin{definition}[Twisted Chern character]
  \label{def:TwistedChernChar}
  \lean{HitchinPW.TwistedChernChar}
  \uses{def:TwistedUniversalFamily}
  \textit{Source: de~Cataldo--Maulik--Shen, Section~0.3.}
  For a twisted universal family \((\mathcal{U}^\alpha, \theta)\), the
  \emph{twisted Chern character} is
  \[
    \ch^\alpha(\mathcal{U}) = \ch(\mathcal{U}) \cup \exp(\alpha)
      \in H^*(C \times \MDol, \QQ),
  \]
  with degree-\(k\) part
  \(\ch_k^\alpha(\mathcal{U}) \in H^{2k}(C \times \MDol, \QQ)\).
  The normalized twisted Chern character is uniquely determined by requiring
  \(\ch_1^\alpha(\mathcal{U})|_{p \times \MDol} = 0\) and
  \(\ch_1^\alpha(\mathcal{U})|_{C \times q} = 0\) for points
  \(p \in \MDol\), \(q \in C\).
\end{definition}

% SOURCE: de~Cataldo--Maulik--Shen, Section~0.3, equation~(0.4)
%   (read from references/MR4433080-hitchin-pw.tex)
% SOURCE QUOTE: For any $\gamma \in H^i(C, \BQ)$, let $c(\gamma ,k)$ denote
% the \emph{tautological class}
% $c(\gamma ,k) := \int_\gamma \mathrm{ch}_k^\alpha(\CU) = {p_{\CM}}_\ast(
%   p_C^\ast \gamma \cup \mathrm{ch}^\alpha_k(\CU)) \in
%   H^{i+2k-2}(\CM_{\mathrm{Dol}}, \BQ).$
\begin{definition}[Tautological class]
  \label{def:TautologicalClass}
  \lean{HitchinPW.TautologicalClass}
  \uses{def:TwistedChernChar}
  \textit{Source: de~Cataldo--Maulik--Shen, Section~0.3, equation~(0.4).}
  For \(\gamma \in H^i(C,\QQ)\), the \emph{tautological class} is
  \[
    c(\gamma,k) := \int_\gamma \ch_k^\alpha(\mathcal{U})
      = (p_{\mathcal{M}})_*(p_C^* \gamma \cup \ch_k^\alpha(\mathcal{U}))
      \in H^{i+2k-2}(\MDol, \QQ).
  \]
  Markman showed that the classes \(c(\gamma,k)\) generate
  \(H^*(\MDol, \QQ)\) as a \(\QQ\)-algebra.
\end{definition}

% SOURCE: de~Cataldo--Maulik--Shen, Section~0.3
%   (read from references/MR4433080-hitchin-pw.tex)
% SOURCE QUOTE: $^k\mathrm{Hdg}^d(\CM_B) = W_{2k}H^d(\CM_B, \BQ) \cap F^k
%   H^d(\CM_B, \BC) \cap \bar{F}^k H^d(\CM_B, \BC),$
% with $F^*$ and $W_*$ the Hodge and weight filtrations, respectively.
\begin{definition}[Hodge--Tate decomposition]
  \label{def:HodgeTateDecomp}
  \lean{HitchinPW.HodgeTateDecomp}
  \uses{def:WeightFiltration}
  \textit{Source: de~Cataldo--Maulik--Shen, Section~0.3.}
  The \emph{Hodge--Tate piece} of degree \((k,d)\) is
  \[
    \HodgeTate{k}{d}(\MB)
      = W_{2k} H^d(\MB,\QQ) \cap F^k H^d(\MB,\mathbb{C})
        \cap \bar{F}^k H^d(\MB,\mathbb{C}),
  \]
  where \(F^*\) is the Hodge filtration and \(W_*\) the weight filtration.
  The rational cohomology \(H^*(\MB,\QQ)\) is of Hodge--Tate type, and
  \[
    H^*(\MB,\QQ) \cong \bigoplus_{k,d} \HodgeTate{k}{d}(\MB).
  \]
  This structure is axiomatized in the Lean formalization.
\end{definition}

% SOURCE: de~Cataldo--Maulik--Shen, Section~0.3, Conjecture~0.2
%   (read from references/MR4433080-hitchin-pw.tex)
% SOURCE QUOTE: \begin{conj}[Equivalent version of P=W]\label{P=W2}
% There exists a splitting $G_*H^*(\CM_{\mathrm{Dol}}, \BQ)$ of the perverse
% filtration associated with the Hitchin fibration (\ref{Hitchin_fib})
% satisfying the following two properties.
% \begin{enumerate}
%   \item[(a)](Tautological classes) Every tautological class $c(\gamma,k)$
%   has perversity $k$ and in fact
%   $c(\gamma,k) \in G_kH^{\ast}(\CM_{\mathrm{Dol}}, \BQ),\quad \forall k
%     \geq 0, \, \forall \gamma \in H^*(C, \BQ).$
%   \item[(b)] (Multiplicativity) The perverse decomposition is
%   multiplicative, i.e.,
%   $\cup: G_k H^d(\CM_\mathrm{Dol}, \BQ)\times G_{k'} H^{d'}(\CM_\mathrm{Dol},
%     \BQ) \to G_{k+k'} H^{d+d'}(\CM_\mathrm{Dol}, \BQ).$
% \end{enumerate}
% \end{conj}
% [restating P=W as $P_k H^* = \bigoplus_{k' \leq k} {}^{k'}\mathrm{Hdg}^*$]
\begin{conjecture}[Equivalent version of P=W]
  \label{conj:PW_splitting}
  \uses{def:PervSplitting, def:HitchinFibration, def:TautologicalClass}
  \textit{Source: de~Cataldo--Maulik--Shen, Section~0.3, Conjecture~0.2.}
  The P=W Conjecture (equivalently, the prediction
  \(P_k H^*(\MDol,\QQ) = \bigoplus_{k' \leq k} \HodgeTate{k'}{*}(\MB)\)) is
  equivalent to the following statement.  There exists a splitting
  \(G_* H^*(\MDol,\QQ)\) of the perverse filtration associated with the
  Hitchin fibration (\cref{def:HitchinFibration}) satisfying the two
  properties:
  \begin{enumerate}
    \item[(a)] \emph{(Tautological classes)} Every tautological class
      \(c(\gamma,k)\) has perversity \(k\); in fact
      \[
        c(\gamma,k) \in G_k H^*(\MDol,\QQ)
        \quad \forall\, k \geq 0, \ \forall\, \gamma \in H^*(C,\QQ).
      \]
    \item[(b)] \emph{(Multiplicativity)} The perverse decomposition is
      multiplicative:
      \[
        \cup \colon G_k H^d(\MDol,\QQ) \times G_{k'} H^{d'}(\MDol,\QQ)
          \longrightarrow G_{k+k'} H^{d+d'}(\MDol,\QQ).
      \]
  \end{enumerate}
\end{conjecture}

% SOURCE: de~Cataldo--Maulik--Shen, Section~0.3, equation~(0.5)
%   (read from references/MR4433080-hitchin-pw.tex)
% SOURCE QUOTE: Furthermore, Shende proved in \cite{Shende}
% { (cf. Lemma \ref{lem3.1})} that
% $c(\gamma,k) \in {^k\mathrm{Hdg}^{i+2k-2}(\CM_B)},
%   \quad \forall~\gamma \in H^i(C, \BQ),$
% where $^k\mathrm{Hdg}^d(\CM_B) = W_{2k}H^d(\CM_B, \BQ) \cap F^k
%   H^d(\CM_B, \BC) \cap \bar{F}^k H^d(\CM_B, \BC),$ with $F^*$ and $W_*$
% the Hodge and weight filtrations, respectively.
\begin{lemma}[Tautological weight]
  \label{lem:taut_weight}
  \lean{HitchinPW.taut_weight}
  \uses{def:TautologicalClass, def:HodgeTateDecomp}
  \textit{Source: de~Cataldo--Maulik--Shen, equation~(0.5), citing Shende.}
  For all \(\gamma \in H^i(C,\QQ)\) and \(k \geq 0\),
  \[
    c(\gamma,k) \in \HodgeTate{k}{i+2k-2}(\MB).
  \]
  In particular, \(c(\gamma,k) \in W_{2k} H^{i+2k-2}(\MB,\QQ)\).
\end{lemma}

\begin{proof}
  \uses{}
  This is proved by Shende (see Lemma~3.1 of the paper).  The argument uses
  the explicit description of the mixed Hodge structure on \(\MB\) and the
  fact that the tautological classes arise from the universal family over
  \(C \times \MDol\) via cup product and pushforward, operations that are
  compatible with the Hodge filtration.  The proof is axiomatized in the
  Lean formalization.
\end{proof}

\subsection{P=W for tautological classes}

% SOURCE: de~Cataldo--Maulik--Shen, Section~0.4.1
%   (read from references/MR4433080-hitchin-pw.tex)
% SOURCE QUOTE: We consider the $\BQ$-subalgebra
% $R^*(\CM_{\mathrm{Dol}})  \subset H^*(\CM_{\mathrm{Dol}}, \BQ)=
%   H^*(\CM_B, \BQ)$
% generated by all the \emph{even} tautological classes
% $c(\gamma,k), \quad  k\geq 0, \quad  \gamma \in H^0(C, \BQ)\oplus
%   H^2(C, \BQ).$
\begin{definition}[Even tautological subalgebra]
  \label{def:EvenTautSubalgebra}
  \lean{HitchinPW.EvenTautSubalgebra}
  \uses{def:TautologicalClass}
  \textit{Source: de~Cataldo--Maulik--Shen, Section~0.4.1.}
  The \emph{even tautological subalgebra} is the \(\QQ\)-subalgebra
  \[
    R^*(\MDol) \subset H^*(\MDol, \QQ)
  \]
  generated by all even tautological classes
  \(c(\gamma,k)\) with \(k \geq 0\) and
  \(\gamma \in H^0(C,\QQ) \oplus H^2(C,\QQ)\).
\end{definition}

% SOURCE: de~Cataldo--Maulik--Shen, Section~0.4.1, Theorem~0.2
%   (read from references/MR4433080-hitchin-pw.tex)
% SOURCE QUOTE: \begin{thm}\label{P=W_even}
% The P=W conjecture holds, for any genus $g \geq 2$ and rank $r\geq 1$ for
% $R^*(\CM_{\mathrm{Dol}})$, \emph{i.e.}
% $P_kH^\ast(\CM_{\mathrm{Dol}}, \BQ) \cap R^*(\CM_{\mathrm{Dol}})
%   = W_{2k}H^\ast(\CM_{B}, \BQ) \cap R^*(\CM_{\mathrm{Dol}})
%   = W_{2k+1}H^\ast(\CM_{B}, \BQ) \cap R^*(\CM_{\mathrm{Dol}}).$
% \end{thm}
\begin{theorem}[P=W for even tautological classes]
  \label{thm:PW_even}
  \lean{HitchinPW.PW_even}
  \uses{def:EvenTautSubalgebra, def:PervFiltration, def:WeightFiltration}
  \textit{Source: de~Cataldo--Maulik--Shen, Theorem~0.2.}
  For any \(g \geq 2\) and \(r \geq 1\), the P=W equality holds on the even
  tautological subalgebra \(R^*(\MDol)\):
  \[
    P_k H^*(\MDol,\QQ) \cap R^*(\MDol)
      = W_{2k} H^*(\MB,\QQ) \cap R^*(\MDol)
      = W_{2k+1} H^*(\MB,\QQ) \cap R^*(\MDol)
    \quad \forall\, k \geq 0.
  \]
\end{theorem}

% SOURCE QUOTE PROOF: [The proof in Section 4 uses the abelian surface
% splitting Theorem \ref{taut_abelian}, the specialization map
% compatibility, and Markman's monodromy operators to reduce to the compact
% case.]
\begin{proof}
  \uses{thm:taut_abelian_splitting, lem:taut_weight}
  The proof proceeds by a degeneration argument.  One embeds a genus-\(g\)
  curve \(C\) into an abelian surface \(A\) and considers the degeneration
  to the normal cone.  The specialization map on cohomology is shown to be
  compatible with tautological classes and perverse splittings.  Applying
  the abelian surface splitting (\cref{thm:taut_abelian_splitting}) and
  the weight-filtration identity \cref{lem:taut_weight} to the even
  tautological generators of \(R^*(\MDol)\) yields the equality
  \(P_k \cap R^* = W_{2k} \cap R^*\).  The multiplicativity of the abelian
  splitting (\cref{def:PervSplitting}) propagates the equality from
  generators to the full subalgebra.
  The specialization argument uses Markman's monodromy operators, which are
  axiomatized in the Lean formalization; the full argument is in Section~4
  of the paper.
\end{proof}

% SOURCE: de~Cataldo--Maulik--Shen, Section~0.4.2, Theorem~0.3
%   (read from references/MR4433080-hitchin-pw.tex)
% SOURCE QUOTE: \begin{thm}\label{P=W_odd}
% Any odd tautological class (\ref{odd_taut}) has perversity $k$.
% \end{thm}
% [with odd tautological classes defined as:]
% $c(\gamma, k) \in H^{2k-1}(\CM_{\mathrm{Dol}}, \BQ), \quad \forall k
%   \geq 0, \, \forall \gamma \in H^1(C, \BQ).$
\begin{theorem}[Odd tautological class perversity]
  \label{thm:oddTautPerversity}
  \lean{HitchinPW.oddTautPerversity}
  \uses{def:TautologicalClass, def:PervFiltration}
  \textit{Source: de~Cataldo--Maulik--Shen, Theorem~0.3.}
  For any \(\gamma \in H^1(C,\QQ)\) and \(k \geq 0\), the odd tautological
  class
  \[
    c(\gamma,k) \in H^{2k-1}(\MDol,\QQ)
  \]
  has perversity exactly \(k\): that is,
  \(c(\gamma,k) \in P_k H^*(\MDol,\QQ)\) but
  \(c(\gamma,k) \notin P_{k-1} H^*(\MDol,\QQ)\).
\end{theorem}

% SOURCE QUOTE PROOF: [The proof uses the Curious Hard Lefschetz theorem
% (Mellit's theorem, axiomatized) together with the abelian surface results
% to show the perversity bound, and the weight identity
% (\ref{taut_weight}) to establish the exact perversity.]
\begin{proof}
  \uses{thm:taut_abelian_splitting, lem:taut_weight}
  The upper bound \(c(\gamma,k) \in P_k\) follows from the abelian surface
  specialization and \cref{lem:taut_weight}, which shows
  \(c(\gamma,k) \in \HodgeTate{k}{2k-1}(\MB)\) and hence
  \(c(\gamma,k) \in W_{2k} H^{2k-1}(\MB,\QQ)\).  The weight filtration
  inclusion \(W_{2k} \subset P_k\) (a consequence of P=W for even classes
  via \cref{thm:PW_even}) gives the upper bound.  The strict lower bound
  \(c(\gamma,k) \notin P_{k-1}\) is proved using the Curious Hard Lefschetz
  theorem due to Mellit, which is axiomatized in the Lean formalization;
  the full argument is in Section~4 of the paper.
\end{proof}

% SOURCE: de~Cataldo--Maulik--Shen, Section~0.4.3, Theorem~0.4
%   (read from references/MR4433080-hitchin-pw.tex)
% SOURCE QUOTE: \begin{thm}\label{last}
% The P=W Conjecture \ref{P=W_conj} is equivalent to the multiplicativity
% (\ref{multi1}) of the perverse filtration.
% \end{thm}
% [multiplicativity defined as:]
% $\cup: P_k H^d(\CM_\mathrm{Dol}, \BQ)\times P_{k'} H^{d'}(\CM_\mathrm{Dol},
%   \BQ) \to P_{k+k'} H^{d+d'}(\CM_\mathrm{Dol}, \BQ).$
\begin{theorem}[Equivalence with multiplicativity]
  \label{thm:PW_iff_multiplicativity}
  \lean{HitchinPW.PW_iff_multiplicativity}
  \uses{thm:PW_even, thm:oddTautPerversity, def:PervFiltration}
  \textit{Source: de~Cataldo--Maulik--Shen, Theorem~0.4.}
  The P=W Conjecture is equivalent to the multiplicativity of the perverse
  filtration:
  \[
    \cup \colon P_k H^d(\MDol,\QQ) \times P_{k'} H^{d'}(\MDol,\QQ)
      \longrightarrow P_{k+k'} H^{d+d'}(\MDol,\QQ) \quad \forall\, k, k', d, d'.
  \]
\end{theorem}

% SOURCE QUOTE PROOF: In fact, we deduce from (\ref{taut_weight}),
% (\ref{eqn9}), and (\ref{multi1}) that
% $W_{2i}H^*(\CM_B, \BQ) \subset P_iH^*(\CM_{\mathrm{Dol}}, \BQ),
%   \quad \forall i\geq 0,$
% which further implies Conjecture \ref{P=W_conj} by
% \cite[Theorem 1.5.3]{Mellit} and the second paragraph following
% \cite[Corollary 1.5.4]{Mellit}.
\begin{proof}
  \uses{thm:PW_even, thm:oddTautPerversity, lem:taut_weight}
  The implication P=W \(\Rightarrow\) multiplicativity is immediate from the
  weight filtration being multiplicative (standard mixed Hodge theory).  For
  the converse: \cref{lem:taut_weight} and \cref{thm:oddTautPerversity}
  together show that all tautological generators \(c(\gamma,k)\) satisfy
  \[
    c(\gamma,k) \in \HodgeTate{k}{*}(\MB) \subset W_{2k} H^*(\MB,\QQ)
      \subset P_k H^*(\MDol,\QQ).
  \]
  Assuming multiplicativity of \(P_*\) and using that the \(c(\gamma,k)\)
  generate all of \(H^*(\MDol,\QQ)\), one deduces
  \(W_{2i} H^*(\MB,\QQ) \subset P_i H^*(\MDol,\QQ)\).
  The full P=W equality then follows from Mellit's Theorem~1.5.3
  (axiomatized in the Lean formalization).
\end{proof}

\section{Perverse filtrations}

\subsection{Perverse sheaves}

% ============================================================
%  Definitions 1--2: Perverse filtrations
% ============================================================

% SOURCE: de~Cataldo--Maulik--Shen, Section~1.2, equation~(1.2)
%   (read from references/MR4433080-hitchin-pw.tex)
% SOURCE QUOTE: For $k \in \BZ$, let $^\mathfrak{p}\tau_{\leq k}$ be the
% truncation functor associated with the perverse $t$-structure. Given an
% object $\CC \in D_c^b(Y)$, there is a natural morphism
% $^\mathfrak{p}\tau_{\le k}\CC \rightarrow \CC.$
% For the morphism $\pi:X \to Y$, we obtain from (\ref{trun0}) the morphism
% $^\mathfrak{p}\tau_{\le k}R\pi_\ast \BQ_X \rightarrow R\pi_\ast \BQ_X,$
% which further induces a morphism of (hyper-)cohomology groups
% $H^{d-(\dim X -{ R})}\Big{(}Y, \, ^\mathfrak{p}\tau_{\le k}(R\pi_\ast
%   \BQ_X[\dim X - { R}]) \Big{)} \rightarrow H^d(X, \BQ).$
% Here ${R} = \dim X \times_Y X - \dim X$ is the \emph{defect of
% semismallness}.  The $k$-th piece of the perverse filtration
% $P_kH^d(X, \BQ) \subset H^d(X, \BQ)$
% is defined to be the image of (\ref{perv_filt}).
\begin{definition}[Perverse filtration]
  \label{def:PervFiltration}
  \lean{HitchinPW.PerverseFiltration}
  \uses{}
  \textit{Source: de~Cataldo--Maulik--Shen, Section~1.2, equation~(1.2).}
  Let \(\pi \colon X \to Y\) be a proper morphism with \(X\) a nonsingular
  algebraic variety over \(\mathbb{C}\).  Write
  \(R = \dim(X \times_Y X) - \dim X\) for the \emph{defect of semismallness}.
  The \emph{perverse filtration} associated with \(\pi\) is the increasing
  filtration
  \[
    P_0 H^*(X,\QQ) \subset P_1 H^*(X,\QQ) \subset \cdots \subset H^*(X,\QQ),
  \]
  where the \(k\)-th piece \(P_k H^d(X,\QQ)\) is the image of the natural map
  \[
    H^{d-(\dim X - R)}\!\bigl(Y,\;
      {}^{\mathfrak{p}}\tau_{\leq k}(R\pi_* \QQ_X[\dim X - R])\bigr)
    \longrightarrow H^d(X,\QQ)
  \]
  induced by the perverse truncation morphism
  \({}^{\mathfrak{p}}\tau_{\leq k} R\pi_* \QQ_X \to R\pi_* \QQ_X\).
  This filtration is axiomatized in the Lean formalization.
\end{definition}

% SOURCE: de~Cataldo--Maulik--Shen, Section~1.2
%   (read from references/MR4433080-hitchin-pw.tex)
% SOURCE QUOTE: We say that a graded vector space decomposition
% $H^*(X, \BQ) = \bigoplus_{k,d} G_k H^d(X, \BQ)$
% \emph{splits} the perverse filtration, if
% $P_k H^d(X, \BQ) = \bigoplus_{i \leq k} G_i H^d(X, \BQ), \quad \forall k,d.$
\begin{definition}[Perverse splitting]
  \label{def:PervSplitting}
  \lean{HitchinPW.PerverseSplitting}
  \uses{def:PervFiltration}
  \textit{Source: de~Cataldo--Maulik--Shen, Section~1.2.}
  A graded vector space decomposition
  \(H^*(X,\QQ) = \bigoplus_{k,d} G_k H^d(X,\QQ)\)
  \emph{splits the perverse filtration} associated with \(\pi\) if
  \[
    P_k H^d(X,\QQ) = \bigoplus_{i \leq k} G_i H^d(X,\QQ) \quad \forall\, k, d.
  \]
\end{definition}

\subsection{Perverse filtration for projective bases}

% SOURCE: de~Cataldo--Maulik--Shen, Section~1.4, Proposition~1.1
%   (read from references/MR4433080-hitchin-pw.tex)
% SOURCE QUOTE: \begin{prop}[\cite{dCM0} Proposition 5.2.4.] \label{prop1.1}
% Assume that
% $\dim X = 2 \dim Y = 2R.$
% Then we have
% $P_kH^m(X, \BQ) = \sum_{i\geq 1} \left(
%   \mathrm{Ker}(L^{R+k+i-m}) \cap \mathrm{Im}(L^{i-1}) \right)
%   \cap H^m(X, \BQ).$
% \end{prop}
\begin{proposition}[Perverse filtration via ample class]
  \label{prop:perv_ample}
  \lean{HitchinPW.perv_ample}
  \uses{def:PervFiltration}
  \textit{Source: de~Cataldo--Maulik--Shen, Proposition~1.1 (citing
  de~Cataldo--Migliorini, Proposition~5.2.4).}
  Let \(\pi \colon X \to Y\) be a proper morphism with
  \(\dim X = 2\dim Y = 2R\), and let \(\eta\) be an ample class on \(Y\).
  Set \(L = \pi^* \eta \in H^2(X,\QQ)\).  Then
  \[
    P_k H^m(X,\QQ)
      = \sum_{i \geq 1}
        \bigl(\ker(L^{R+k+i-m}) \cap \operatorname{Im}(L^{i-1})\bigr)
        \cap H^m(X,\QQ).
  \]
\end{proposition}

% SOURCE QUOTE PROOF: [Cited from de Cataldo--Migliorini \cite{dCM0}
% Proposition 5.2.4; the proof uses the decomposition theorem and Relative
% Hard Lefschetz.]
\begin{proof}
  \uses{}
  This is cited from de~Cataldo--Migliorini, Proposition~5.2.4.  The proof
  uses the Decomposition Theorem and the Relative Hard Lefschetz theorem
  applied to the morphism \(\pi\).  The result is axiomatized in the Lean
  formalization.
\end{proof}

\subsection{Abelian surfaces and Hilbert schemes}

% SOURCE: de~Cataldo--Maulik--Shen, Section~1.5, Theorem~1.6
%   (read from references/MR4433080-hitchin-pw.tex)
% SOURCE QUOTE: \begin{thm}\label{taut_Hilb}
% Let $\CI_n$ be the universal ideal sheaf on $A \times A^{[n]}$.
% \begin{enumerate}
%   \item[(a)](Tautological classes) We have
%   $\mathrm{ch}_k(\CI_n) \in \widetilde{G}_k H^{2k}(A\times A^{[n]}, \BQ).$
%   \item[(b)](Multiplicativity) The decomposition  (\ref{canonical_split})
%   is multiplicative, i.e.
%   $\cup: \widetilde{G}_i H^d(A^{[n]}, \BQ) \times \widetilde{G}_{i'}
%     H^{d'}(A^{[n]}, \BQ) \to \widetilde{G}_{i+i'}H^{d+d'}(A^{[n]}, \BQ).$
% \end{enumerate}
% \end{thm}
\begin{theorem}[Hilbert scheme case]
  \label{thm:taut_Hilb}
  \lean{HitchinPW.taut_Hilb}
  \uses{def:PervSplitting}
  \textit{Source: de~Cataldo--Maulik--Shen, Theorem~1.6.}
  Let \(A = E \times E'\) be a product of elliptic curves, and let
  \(p_n \colon A^{[n]} \to (E')^{(n)}\) be the natural morphism from the
  Hilbert scheme of \(n\) points.  Let \(\mathcal{I}_n\) be the universal
  ideal sheaf on \(A \times A^{[n]}\).  The canonical splitting
  \(\widetilde{G}_* H^*(A^{[n]},\QQ)\) of the perverse filtration
  associated with \(p_n\) satisfies:
  \begin{enumerate}
    \item[(a)] \emph{(Tautological classes)}
      \[
        \ch_k(\mathcal{I}_n) \in \widetilde{G}_k H^{2k}(A \times A^{[n]},\QQ).
      \]
    \item[(b)] \emph{(Multiplicativity)} The decomposition is multiplicative:
      \[
        \cup \colon \widetilde{G}_i H^d(A^{[n]},\QQ)
          \times \widetilde{G}_{i'} H^{d'}(A^{[n]},\QQ)
          \longrightarrow \widetilde{G}_{i+i'} H^{d+d'}(A^{[n]},\QQ).
      \]
  \end{enumerate}
\end{theorem}

% SOURCE QUOTE PROOF: Theorem \ref{taut_Hilb} (b) was proven in \cite{Z}.
% Although the multiplicativity was shown only for the perverse filtration
% in \cite{Z}, the same proof works for the decomposition
% (\ref{canonical_split}).  This was explained in
% \cite[Proposition 1.8]{SZ}.  Theorem \ref{taut_Hilb} (a) follows directly
% from the proof of \cite[Theorem 0.1]{SZ} and the proof of
% \cite[Theorem 3.3]{SZ}, where it is explained how to treat perverse
% decompositions instead of perverse filtrations.
\begin{proof}
  \uses{}
  Part~(a) follows from the proof of Zhang--Shen (SZ, Theorems 0.1 and 3.3),
  using the explicit formula~(1.10) of the paper for the canonical splitting
  and the Nakajima operators on Hilbert schemes.
  Part~(b) (multiplicativity) was proved by Zhang (Z) for the perverse
  filtration; the same argument applies to the canonical splitting decomposition,
  as explained in Shen--Zhang (SZ, Proposition~1.8).
  Both results are axiomatized in the Lean formalization via sorry-proofs.
\end{proof}

\section{Moduli of one-dimensional sheaves}

\subsection{Overview and main results}

% ============================================================
%  Definitions 13--14: Abelian surface moduli
% ============================================================

% SOURCE: de~Cataldo--Maulik--Shen, Section~2.1
%   (read from references/MR4433080-hitchin-pw.tex)
% SOURCE QUOTE: Throughout this section, we assume $A$ is an abelian surface
% and fix $\chi \in \mathbb{Z}$.  Let $\beta \in H_2(A, \BZ)$ be an ample
% curve class with $\beta^2 \geq 6$; we only consider classes $\beta$ for
% which the vector $(0,\beta,\chi)\in H^{\mathrm{ev}}(A,\BZ)$ is primitive.
% ...
% Let $\CM_{\beta,A}$ be the moduli space which parametrizes one-dimensional
% sheaves satisfying
% $[\mathrm{supp}(\CF)]= \beta, \quad \chi(\CF) = \chi,$
% that are Gieseker-stable with respect to our generic polarization.
\begin{definition}[Moduli of sheaves on an abelian surface]
  \label{def:ModuliSheaves}
  \lean{HitchinPW.ModuliSheaves}
  \uses{}
  \textit{Source: de~Cataldo--Maulik--Shen, Section~2.1.}
  Let \(A\) be an abelian surface, \(\beta \in H_2(A,\ZZ)\) an ample curve
  class with \(\beta^2 \geq 6\) and primitive Mukai vector
  \((0,\beta,\chi) \in H^{\mathrm{ev}}(A,\ZZ)\), and let \(H\) be a
  polarization generic with respect to \((0,\beta,\chi)\).  The moduli space
  \(\mathcal{M}_{\beta,A}\) parametrizes Gieseker \(H\)-stable one-dimensional
  sheaves \(\mathcal{F}\) on \(A\) satisfying
  \[
    [\supp(\mathcal{F})] = \beta, \quad \chi(\mathcal{F}) = \chi.
  \]
  It is a smooth projective variety of dimension \(\beta^2 + 2\).
\end{definition}

% SOURCE: de~Cataldo--Maulik--Shen, Section~2.1, equation~(2.3)
%   (read from references/MR4433080-hitchin-pw.tex)
% SOURCE QUOTE: The moduli space $\CM_{\beta,A}$ admits a natural
% Hilbert--Chow morphism
% $\pi_\beta: \CM_{\beta,A} \to B, \quad \pi_\beta(\CF) = [\mathrm{supp}(\CF)]$
% where $B$ is the Chow variety parametrizing effective one-cycles in the
% class $\beta$; see \cite{LP1, LP2}.
\begin{definition}[Hilbert--Chow morphism]
  \label{def:HilbChowMorphism}
  \lean{HitchinPW.HilbChowMorphism}
  \uses{def:ModuliSheaves}
  \textit{Source: de~Cataldo--Maulik--Shen, Section~2.1, equation~(2.3).}
  The \emph{Hilbert--Chow morphism} is the natural map
  \[
    \pi_\beta \colon \mathcal{M}_{\beta,A} \longrightarrow B,
    \quad \mathcal{F} \mapsto [\supp(\mathcal{F})],
  \]
  where \(B\) is the Chow variety of effective one-cycles in the class
  \(\beta\).
\end{definition}

% ============================================================
%  Main Theorem 15: Splitting for abelian surface moduli
% ============================================================

% SOURCE: de~Cataldo--Maulik--Shen, Section~2.1, Theorem~2.1
%   (read from references/MR4433080-hitchin-pw.tex)
% SOURCE QUOTE: \begin{thm}\label{taut_abelian}
% There exists a choice of universal twisted class
% $\mathrm{ch}^\alpha(\CF_\beta)$ with $\alpha$ of type (\ref{alpha}),
% and a splitting $\widetilde{G}_\ast H^\ast(\CM_{\beta,A}, \BC)$ of the
% perverse filtration associated with the Hilbert-Chow morphism $\pi_\beta$
% (\ref{Hilb-Ch}), satisfying the following properties.
% \begin{enumerate}
%   \item[(a)](Tautological classes) For any $\gamma \in H^*(A, \BC)$, we have
%   $\int_\gamma \mathrm{ch}^\alpha_k(\CF_\beta) \in
%     \widetilde{G}_{k-1}H^*(\CM_{\beta,A}, \BC).$
%   \item[(b)](Multiplicativity) We have
%   $\cup: \widetilde{G}_i H^d(\CM_{\beta,A}, \BC) \times
%     \widetilde{G}_{i'} H^{d'}(\CM_{\beta,A}, \BC) \to
%     \widetilde{G}_{i+i'}H^{d+d'}(\CM_{\beta,A}, \BC).$
% \end{enumerate}
% \end{thm}
\begin{theorem}[Splitting for abelian surface moduli]
  \label{thm:taut_abelian_splitting}
  \lean{HitchinPW.taut_abelian_splitting}
  \uses{def:ModuliSheaves, def:HilbChowMorphism, def:PervSplitting,
        def:TwistedChernChar, def:TautologicalClass}
  \textit{Source: de~Cataldo--Maulik--Shen, Theorem~2.1.}
  There exist a choice of universal twisted class
  \(\ch^\alpha(\mathcal{F}_\beta)\) and a splitting
  \(\widetilde{G}_* H^*(\mathcal{M}_{\beta,A}, \mathbb{C})\) of the perverse
  filtration associated with the Hilbert--Chow morphism
  \(\pi_\beta\) (\cref{def:HilbChowMorphism}) such that:
  \begin{enumerate}
    \item[(a)] \emph{(Tautological classes)} For any
      \(\gamma \in H^*(A,\mathbb{C})\),
      \[
        \int_\gamma \ch_k^\alpha(\mathcal{F}_\beta)
          \in \widetilde{G}_{k-1} H^*(\mathcal{M}_{\beta,A}, \mathbb{C}).
      \]
    \item[(b)] \emph{(Multiplicativity)} The splitting is multiplicative:
      \[
        \cup \colon \widetilde{G}_i H^d(\mathcal{M}_{\beta,A}, \mathbb{C})
          \times \widetilde{G}_{i'} H^{d'}(\mathcal{M}_{\beta,A}, \mathbb{C})
          \longrightarrow
          \widetilde{G}_{i+i'} H^{d+d'}(\mathcal{M}_{\beta,A}, \mathbb{C}).
      \]
  \end{enumerate}
\end{theorem}

% SOURCE QUOTE PROOF: Theorem \ref{taut_abelian} will be deduced from
% Theorem \ref{taut_Hilb} and Markman's results on the monodromy of
% holomorphic symplectic varieties \cite{Markman2, Markman3}.
% The strategy of the argument is to use Theorem \ref{taut_Hilb} and direct
% computations to study the case of primitive curve classes for special
% abelian varieties.  To pass to imprimitive curve classes (which form a
% distinct deformation class of Lagrangian fibrations), we require a mild
% extension of Markman's work to the setting of complex coefficients.
\begin{proof}
  \uses{thm:taut_Hilb}
  The result is deduced from \cref{thm:taut_Hilb} (the Hilbert scheme case)
  and Markman's monodromy operators for holomorphic symplectic varieties.
  For primitive classes \(\beta\), the splitting is constructed directly via
  explicit computations using Theorem~1.6 and Markman's monodromy results.
  The primitive case is then extended to imprimitive \(\beta\) using
  Markman's parallel transport operators, which require a mild generalization
  to complex coefficients.  The monodromy infrastructure is axiomatized in
  the Lean formalization; the full argument is in Section~2 of the paper.
\end{proof}
